\chapter{The widget catalogue, by architecture}
\label{ch:widget-arch}

\begin{aside}
\maya{} ships dozens of widgets in \file{maya/include/maya/widget/}.
Each is a small builder that produces an \code{Element}. The
catalogue is too large to enumerate here (that is
Chapter~\ref{ch:widgets}), but the \emph{architecture} of a
widget --- the patterns every widget follows --- is small enough
to learn in one sitting. This chapter is the architectural
template.
\end{aside}

\section{What a widget is, exactly}

A widget is one of three things:

\begin{itemize}[leftmargin=*]
  \item \textbf{A pure renderer.} Takes parameters, returns an
    \code{Element}. No state, no input handling. The
    \code{badge}, \code{divider}, \code{spinner} widgets.
  \item \textbf{A stateful component.} Has its own \code{Model},
    \code{Msg}, \code{update}, \code{view}. Composes via
    \code{Cmd::map}. The \code{textarea}, \code{picker},
    \code{table} widgets.
  \item \textbf{A reactive widget.} Uses signals to drive
    self-contained animations or watched state. The
    \code{spinner} (running variant), \code{progress},
    \code{streaming\_cursor}.
\end{itemize}

The three layer cleanly: a pure renderer can be wrapped by a
stateful component, which can be wrapped by a reactive widget.
The \code{Element} type at the boundary lets them compose.

\section{The pure-renderer pattern}

Smallest possible widget:

\begin{lstlisting}
struct Badge {
    std::string label;
    Color fg = Color::default_();
    Color bg = Color::default_();

    Badge& with_label(std::string s) { label = std::move(s); return *this; }
    Badge& with_fg(Color c)         { fg = c; return *this; }
    Badge& with_bg(Color c)         { bg = c; return *this; }

    Element build() && {
        return text(\"[\" + label + \"]\")
            | Style{.fg = fg, .bg = bg, .attributes = Attribute::Bold};
    }
    operator Element() && { return std::move(*this).build(); }
};

Badge badge(std::string label = \"\") {
    return Badge{.label = std::move(label)};
}
\end{lstlisting}

Three setters, one builder, one conversion. The \code{build}
produces a styled text element. Trivially composable in the DSL:

\begin{lstlisting}
auto bar = h(
    widget(badge(\"NORMAL\").with_bg(Color::rgb(0,80,150))),
    text(\" file.txt\")
);
\end{lstlisting}

\subsection*{Pure renderer characteristics}

\begin{itemize}[leftmargin=*]
  \item Zero allocations beyond the produced element.
  \item No subscription requirements.
  \item Trivially testable: call \code{build}, inspect the result.
  \item No memory between frames; each \code{build} is fresh.
\end{itemize}

Use the pure-renderer pattern for any widget that is just a
visual arrangement. Badges, dividers, callouts, tags, labels,
chips, breadcrumbs.

\section{The stateful-component pattern}

A widget with its own state, message, update:

\begin{lstlisting}
namespace textarea {

struct Model {
    std::vector<std::string> lines;
    size_t cursor_row = 0, cursor_col = 0;
    size_t scroll_top = 0;
    bool focused = false;
};

struct InsertChar { char c; };
struct Backspace {};
struct CursorMove { int dr, dc; };
struct GainFocus {};
struct LoseFocus {};
using Msg = std::variant<InsertChar, Backspace, CursorMove,
                         GainFocus, LoseFocus>;

std::pair<Model, Cmd<Msg>> update(Model m, Msg msg);
Element view(const Model& m, int width, int height);
Sub<Msg> subscribe(const Model& m);

}
\end{lstlisting}

This is the same shape as a \maya{} application! That is the
point: widgets are sub-apps. Composing them is exactly the
parent/child pattern from Chapter~\ref{ch:elm}.

\subsection*{Parent integration}

\begin{lstlisting}
struct App {
    struct Model { textarea::Model editor; };
    struct ToEditor { textarea::Msg m; };
    using Msg = std::variant<ToEditor, Quit>;

    static std::pair<Model, Cmd<Msg>>
    update(Model m, Msg msg) {
        return std::visit(overload{
            [&](ToEditor te) {
                auto [e2, ecmd] = textarea::update(m.editor, te.m);
                m.editor = std::move(e2);
                return std::pair{m,
                    map(ecmd, [](auto x) -> Msg { return ToEditor{x}; })};
            },
            [&](Quit) { return std::pair{m, Cmd<Msg>::quit()}; }
        }, msg);
    }

    static Element view(const Model& m) {
        return textarea::view(m.editor, 80, 24);
    }

    static Sub<Msg> subscribe(const Model& m) {
        return Sub<Msg>::batch({
            Sub<Msg>::map(textarea::subscribe(m.editor),
                          [](auto x) -> Msg { return ToEditor{x}; }),
            // ... app-level subscriptions
        });
    }
};
\end{lstlisting}

The familiar \code{map} threading. Widget messages stay in the
widget's namespace; the parent's \code{Msg} variant has one
alternative per widget instance.

\section{The reactive widget pattern}

A widget that owns a signal for self-driven updates:

\begin{lstlisting}
class Spinner {
public:
    Spinner();
    Element build() &&;

private:
    Signal<int> frame_{0};
    // The signal advances on each animation frame.
};

Spinner::Spinner() {
    // Set up an animation-frame subscriber.
    AnimationFrame{[this] { frame_.set(frame_() + 1); }};
}

Element Spinner::build() && {
    constexpr std::array glyphs = {\"|\", \"/\", \"-\", \"\\\\\"};
    return text(glyphs[frame_() % glyphs.size()]);
}
\end{lstlisting}

The widget owns its signal; the build reads it; the animation
runs automatically. No parent integration needed for the
animation itself.

\subsection*{When reactive widgets shine}

Reactive widgets are best for:

\begin{itemize}[leftmargin=*]
  \item Decorative elements (spinners, pulsing cursors).
  \item Independent live values (clock, frame counter).
  \item Watched filesystem or network state.
\end{itemize}

For business logic, prefer the stateful-component pattern; the
reactive pattern's hidden state hurts testability.

\section{Widget composition}

Widgets compose three ways:

\subsection*{Vertical composition}

A widget contains another widget:

\begin{lstlisting}
class Modal {
    Element content;
    std::string title;

    Element build() && {
        return box(
            v(text(title) | Bold,
              std::move(content)))
            | border<Heavy> | pad<2>;
    }
};
\end{lstlisting}

The outer widget takes the inner as a parameter and embeds it.

\subsection*{Horizontal composition}

Several widgets share a frame:

\begin{lstlisting}
auto bar = h(
    widget(Badge{}.with_label(\"NORMAL\")),
    text(\" \") | grow<1>,
    widget(Clock{})
);
\end{lstlisting}

The DSL composes them; each widget contributes a sub-tree.

\subsection*{Decoration}

A widget wraps another with style:

\begin{lstlisting}
class Panel {
    Element content;
    std::string title;

    Element build() && {
        return box(
            v(text(title) | Bold,
              std::move(content)))
            | border<Rounded> | pad<1>;
    }
};
\end{lstlisting}

Same structural shape as Modal; different decoration. The pattern
recurs.

\section{Theming widgets}

A widget should use roles, not raw colours. The badge example
above takes \code{fg} and \code{bg} as parameters; a
theme-aware variant would default to roles:

\begin{lstlisting}
Element Badge::build() && {
    auto& theme = current_theme();
    Color f = fg.is_default() ? theme.palette[Role::Accent] : fg;
    Color b = bg.is_default() ? theme.palette[Role::Background] : bg;
    return text(\"[\" + label + \"]\")
        | Style{.fg = f, .bg = b, .attributes = Attribute::Bold};
}
\end{lstlisting}

The widget asks for the current theme's accent colour if the
caller did not specify. Theme swaps update the widget for free.

\section{Widget testing}

\subsection*{Pure renderers}

\begin{lstlisting}
TEST_CASE(\"badge renders with brackets\") {
    auto e = badge(\"hi\").with_fg(Color::rgb(0,0,0)).build();
    REQUIRE(text_view(e) == \"[hi]\");
}
\end{lstlisting}

Direct. No setup.

\subsection*{Stateful components}

Same patterns as Chapter~\ref{ch:testing}. The widget's
\code{update} is pure; test it with \code{fold}.

\subsection*{Reactive widgets}

Use \code{test::FakeClock} and a controlled signal driver.
Advance the clock, assert against \code{build()} output.

\section{When to write a widget}

Not every visual snippet needs a widget. A rule:

\begin{itemize}[leftmargin=*]
  \item One-off layout: write it inline.
  \item Reused in three places: extract a function.
  \item Reused in five places, or has state: write a widget.
  \item Public API for other apps: write a widget with docs.
\end{itemize}

Premature widget abstraction costs more than it saves. Wait
until you see the duplication.

\begin{warning}
A widget that takes too many parameters (more than five or
six) is a sign that it is doing two jobs. Split it.
\end{warning}

\section*{Workshop}

\begin{enumerate}[leftmargin=*]
  \item Write a \code{kbd} widget that renders a keyboard
    shortcut (\verb|[Ctrl-X]|) with consistent styling.
  \item Write a \code{toast} stateful component with an
    \code{auto-dismiss} timer (\code{Cmd::after(3s, Dismiss)}).
  \item Read three widgets in \file{maya/include/maya/widget/}
    and classify each as pure, stateful, or reactive.
\end{enumerate}

\begin{recap}
A widget is a pure renderer, a stateful component, or a
reactive widget. Each pattern has a small ceremony; the
\code{Element} type at the boundary lets them compose.
Stateful widgets are sub-apps and integrate via \code{Cmd::map}.
Reactive widgets own signals for self-driven updates. Theme-
awareness is opt-in but cheap. Premature abstraction is the
real enemy.
\end{recap}
